% ------------------------------------------------------------------------
% file `main-bookexercise-1-exercise-body.tex'
%
%     exercise of type `bookexercise' with id `1'
%
% generated by the `exercise' environment of the
%   `xsim' package v0.21 (2022/02/12)
% from source `main' on 2026/10/08 on line 640
% ------------------------------------------------------------------------
1. 关于线性表，下列说法正确的是（　）。\\
A. (银心, 甘尔, 张三) 和 (甘尔, 银心, 张三)是两个相同的线性表 \\
B. 线性表是一种物理结构，线性表的所有元素物理地址一定连续\\
C. 同一个线性表中各个元素的数据类型可以不同\\
D. 不同线性表的数据元素类型可以不同\\

2. 对于顺序表，假设表中已有 \(n\) 个元素，下列操作中时间复杂度一定为 \(O(1)\) 的是（　）。\\
A. 在表中某位置插入一个元素\\
B. 删除表头元素\\
C. 获取第 \(i\) 个位置上的元素\\
D. 查找值为 \(x\) 的元素\\

3. 对于一个带头结点的单链表：
\begin{bulk}
head → A → B → C → NULL
\end{bulk}
若 p 指向节点 A，现在希望在 A 后面插入节点 X，正确的核心操作是（　）。\\
A. p->next\_node = X;
X->next\_node = p->next\_node;\\
B. X->next\_node = p;
p->next\_node = X;\\
C. X->next\_node = p->next\_node;
p->next\_node = X;\\
D. p = X;
X->next\_node = p->next\_node;\\

4. 以下单链表的操作中，时间复杂度一定是O(1)的是（  ）。\\
A. 求表长\\
B. 在表头插入元素\\
C. 在表尾插入元素\\
D. 获取第\(i\)个位置上的元素\\

5. 在\coderef{code:线性表的顺序存储结构}中，我们存储顺序表元素使用的数组是静态存储的，定义其最大存储数量为100.若存满了，则无法再扩充。事实上，
我们可以使用malloc动态分配数据元素的存储空间，而只保存一个该空间的起始地址指针：
\begin{cppcode}
    struct SeqList{
        ElemType* data;         // 数组起始地址
        int current_size;       // 当前表中的元素数量
        // ...
    };
\end{cppcode}
请你想一想，动态存储下，会有哪些操作和静态存储不同？（重点思考一下插入操作）\\

6. 假设数组a中下标0-9处存放了10个int型数据\(a_1\)-\(a_{10}\)，请你用数组a中的数据分别构建以下两个链表，填充对应代码：

1） \(a_1\) → \(a_2\) → ... → \(a_{10}\)
\begin{cppcode}
    LinkList f1(int a[], int n=10) {

    }
\end{cppcode}
2） \(a_{10}\) → \(a_9\) → ... → \(a_1\)
\begin{cppcode}
    LinkList f2(int a[], int n=10) {

    }
\end{cppcode}
建议采用带头结点的实现。\\

7. 请你实现一个算法DeleteValue(LinkList L, ElemType e)，删除带头结点的单链表L中所有值等于e的节点。\\




% 8. 【某年408真题】已知一个带头节点的单链表，节点结构为
% \begin{cppcode}
%     struct ListNode {
%         ElemType data;
%         ListNode* link;
%     };
% \end{cppcode}
% 假设该链表只给出了头指针list。在不改变链表的前提下，请设计一个尽可能高效的算法，查找链表中倒数第k个位置上的节点（k为正整数）。若查找成功，算法输出
% 该节点的data域的值，并返回1；否则，只返回0.要求：\\
% 1）描述算法的基本设计思想；\\
% 2）描述算法的详细实现步骤；\\
% 3）根据设计思想和实现步骤，采用C、C++或Java语言实现算法，关键之处给出简要注释。\\

8. 假设线性表采用带头结点的单链表保存，结点定义如下：
\begin{cppcode}
    typedef struct node
    {
        int data;
        struct node *next;
    }NODE;
\end{cppcode}
请你设计一个尽可能高效的算法，对于任意输入的链表L，都能返回其中间节点。
当表长为奇数时，返回中间的那个节点；当表长为偶数时，返回中间两个节点中左边（位序更小）的那个；当表为空表时，返回空指针NULL。\\

9.

1）请你设计实现将带头结点的单链表逆置的算法。即：将
\begin{bulk}
    \begin{center}
        head → \(a_1\) → \(a_2\) → ... → \(a_n\)
    \end{center}
\end{bulk}
变为
\begin{bulk}
    \begin{center}
        head → \(a_n\) → \(a_{n-1}\) → ... → \(a_1\)
    \end{center}
\end{bulk}

2）设计实现将不带头结点的单链表逆置的算法。\\


10. 【某年408真题】设线性表：
$$ L=(a_1,a_2,a_3,\cdots,a_{n-2},a_{n-1},a_n) $$
采用带头结点的单链表保存，链表中的结点定义如下：
\begin{cppcode}
    typedef struct node
    {
        int data;
        struct node *next;
    }NODE;
\end{cppcode}

请设计一个空间复杂度为 \(O(1)\) 且时间上尽可能高效的算法，重新排列 \(L\) 中的各结点，得到：
$$ L'=(a_1,a_n,a_2,a_{n-1},a_3,a_{n-2},\cdots) $$
要求：\\
1）给出算法的基本设计思想；\\
2）根据设计思想，采用 C 或 C++ 语言描述算法，关键之处给出注释；\\
3）说明所设计算法的时间复杂度。\\
